\documentclass[11pt]{article}
\usepackage[a4paper,margin=25mm]{geometry}
\usepackage{amsmath,amssymb,amsthm,booktabs}
\newtheorem{lemma}{Lemma}
\newcommand{\gen}[1]{\langle#1\rangle}
\title{The finite boundary for two equal blocks}
\author{Bounded proof component for the transitive rank theorem}
\date{12 October 2026}
\begin{document}
\maketitle

We establish the finite obstruction used in the two-block argument.
The classification of the possible groups in this statement is a
separate obligation: no conclusion about other groups is implicit in
the calculation below. In particular, this component is not a proof
of the full transitive rank theorem.

Let $W=S_k\wr C_2$ act on two blocks of size $k$, and let $r$ be the
permutation exchanging their corresponding letters. For a subgroup
$H\leq S_k$ and a character $\chi:H\to C_2$, put
\[
 D_\chi(H)=\{(h,h)r^{\chi(h)}:h\in H\}.
\]
The two symmetric factors in a row $S_a\times S_b$ of the table below
act on disjoint sets of sizes $a,b$ in one block. The two transitive
degree-six actions are the literal permutation groups numbered 9
and 11 in the GAP transitive-group catalogue; their orders are 36
and 48, respectively. The action, rather than the abstract group
alone, is part of the specification.

\begin{lemma}\label{boundary}
Let $H$ be one of the ten actions in Table~1, of degree $k$.
Take an involutory string C-marking $(h_1,\ldots,h_{k-2})$
of $H$, and let $C$ be its image in $D_\chi(H)$. Split this ordered
list as $C_L,C_R$, requiring the two lists to commute elementwise.
There is no tuple
\[
                  (C_L,x,r,y,C_R)
\]
of independent involutions that generates a transitive subgroup of
$W$, satisfies the string commutation relations, has both $x$ and
$y$ noncommuting with $r$, and has $r$ at least as close to an
endpoint as every other generator exchanging the blocks.
\end{lemma}

\begin{proof}
We give the complete finite enumeration and its exhaustion argument.
The accompanying program uses exact finite permutation groups.

First enumerate the involutory string C-markings of $H$. At a prefix
$(h_1,\ldots,h_j)$, an eligible next generator is an involution of $H$
outside the prefix subgroup and commuting with $h_1,\ldots,h_{j-1}$.
We discard a prefix if one of its generators belongs to the subgroup
generated by the others. Every subfamily of an independent family
is independent, so this discards no desired tuple. Also, if the
current subgroup has index less than $2^{k-2-j}$ in $H$, the prefix
cannot acquire the required number of independent generators: each
strict subgroup enlargement has index at least two.

Simultaneous conjugation by $N_{S_k}(H)$ preserves every condition.
At each step choose one orbit representative under the pointwise
stabilizer of the prefix under conjugation, then replace that
stabilizer by its centralizer of the chosen generator. Induction on
the prefix length shows that every ordered marking is represented.
At length $k-2$, require generation of $H$ and test the intersection
property for every pair of subsets of the marking. The test forms
the actual generated subgroups and their intersections, without a
quotient or an assumed presentation.

For each surviving marking enumerate all $2^{k-2}$ assignments of
its generators to $C_2$. Retain exactly those assignments for which
\texttt{GroupHomomorphismByImages} constructs a homomorphism.
Since the marking generates $H$, these are exactly its characters.
Construct the corresponding diagonal lifts explicitly as
permutations of $2k$ letters. Try every cut between the lifts,
including the two empty-end cuts, and require elementwise
commutation across it. Reject a cut and character only if an
already fixed block-exchanging generator is closer to an endpoint
than $r$.

The string relations require $x$ to centralize all of $C_R$ and
all but the last member of $C_L$. Similarly, $y$ must centralize
all of $C_L$ and all but the first member of $C_R$. Compute these
two centralizers inside $W$, enumerate their involutions, and impose
noncommutation with $r$ and the same endpoint restriction. We may
choose one representative for $x$ under the simultaneous centralizer
of $C$ and $r$ in $W$: conjugation by this group preserves the complete
candidate set for $y$. For each remaining pair require $xy=yx$ and
transitivity. Finally test independence by deleting each generator
in turn and comparing the resulting subgroup with the full group.

The exhaustive counts are in Table~1. No transitive completion occurs
in the first eight cases. In the last two cases, all transitive
completions fail independence. Thus no test of the full intersection
property is needed to reject a completion, proving the lemma.
\end{proof}

\begin{table}[ht]
\centering
\begin{tabular}{lrrrr}
\toprule
Action of $H$ & $k$ & Markings & Transitive completions & Independent\\
\midrule
$S_2\times S_3$ &5&4&0&0\\
$S_2\times S_4$ &6&12&0&0\\
$S_3\times S_3$ &6&1&0&0\\
$S_2\times S_5$ &7&4&0&0\\
$S_3\times S_4$ &7&6&0&0\\
$S_2\times S_6$ &8&8&0&0\\
$S_3\times S_5$ &8&2&0&0\\
$S_4\times S_4$ &8&9&0&0\\
$\mathrm{TransitiveGroup}(6,9)$ &6&1&2&0\\
$\mathrm{TransitiveGroup}(6,11)$ &6&12&152&0\\
\bottomrule
\end{tabular}
\caption{Complete bounded search. Markings are represented up to the
normalizer conjugations described in the proof. Completion counts
use the stated pruning of $x$; they are not counts of isomorphism classes.}
\end{table}

The implementation is in \texttt{search.g}; \texttt{run.g} executes
all ten cases, and \texttt{check.py} requires their exact summaries
and successful completion. To check that the procedure can retain
valid completions, it is also run with $k=3$ and the empty marking
of the trivial group. It recovers six independent transitive
rank-three string C-markings after the same pruning. For example,
\[
       ((5\ 6),\ (1\ 4)(2\ 5)(3\ 6),\ (1\ 2))
\]
passes the full intersection test. A redundant triple of distinct
transpositions in $S_3$ supplies negative controls for independence
and the intersection property.

The proof depends on GAP's exact finite permutation-group algorithms
and its degree-six transitive catalogue. It does not supply a Lean
formalization. The separate reduction of an arbitrary two-block
group to these ten actions, the endpoint cases, and all other
branches of the transitive rank theorem remain outside this lemma.
Normalizing the commuting generators by powers of the block swap
already appears in Cameron, Fernandes, Leemans and Mixer,
\emph{String C-groups as transitive subgroups of $S_n$},
arXiv:1410.5863v1 (2014), proof of Proposition~2.1, Case~3.
We retain that step. The present contribution is the explicit
character description and exhaustive finite analysis with two
noncommuting neighbors, within the common string-separator argument.
It neither claims a new open-conjecture solution nor cites the
target rank theorem as a premise.
\end{document}
